\documentclass[11pt,a4paper]{amsart}
\usepackage{amsmath,amssymb,amsthm}
\usepackage{hyperref}

\newtheorem{theorem}{Theorem}[section]
\newtheorem{proposition}[theorem]{Proposition}
\newtheorem{lemma}[theorem]{Lemma}
\newtheorem{corollary}[theorem]{Corollary}
\newtheorem{definition}[theorem]{Definition}
\newtheorem{conjecture}[theorem]{Conjecture}
\newtheorem{remark}[theorem]{Remark}

\title{P04 --- Finite Weil Geometry and the Objekt-X Interface}
\author{Objekt-X-Programm}
\date{2026-08-08 (SYN, in Arbeit)}

\begin{document}
\maketitle

\begin{abstract}
We consolidate NEU-259, 260a, 260b (in progress).
For each $a>0$, Suzuki (2026) constructs RH-free a
self-adjoint operator $A_a$, a Hilbert space $\mathcal{H}(T_a)$,
and a family $\overline{\mathscr{D}}_{a,\theta}$ of self-adjoint
first-order extensions whose eigenvalues converge (conjecturally)
to the imaginary parts of the non-trivial zeros.
We identify the canonical RH-free shift $\lambda_{\mathrm{can}}(a)=\lambda_a-1$,
exhibit the gauge freedom $\theta\mapsto\theta+\beta-\alpha$,
and formulate the arithmetic canonicity problem: does
the BC/adelic involution canonically pair the deficiency
lines $\mathcal{N}_{+,a}\leftrightarrow\mathcal{N}_{-,a}$?
\end{abstract}

\section{Operator chain (Suzuki 2026, RH-free)}

\begin{theorem}[Operator chain, NEU-259 $\checkmark$]
For every $a>0$, the following objects exist unconditionally:
\[
  Q_W^a \longrightarrow A_a \longrightarrow T_{a,\lambda} := A_a-\lambda I
  \longrightarrow \mathcal{H}(T_{a,\lambda})
  \longrightarrow \overline{\mathscr{D}}_{a,\theta}.
\]
Specifically:
\begin{enumerate}
\item $Q_W^a$ is semibounded and closable on $L^2(-a,a)$.
\item $A_a = A_a^*$ with discrete spectrum bounded below.
\item For $\lambda < \lambda_a := \inf\sigma(A_a)$:
  $T_{a,\lambda} > 0$ and $\mathcal{H}(T_{a,\lambda})$ is a Hilbert space.
\item $\overline{\mathscr{D}}_{a,\theta} = (\overline{\mathscr{D}}_{a,\theta})^*$
  on $\mathcal{H}(T_{a,\lambda})$ for each $\theta\in[0,2\pi)$,
  deficiency indices $(1,1)$.
\end{enumerate}
\end{theorem}

\section{Canonical RH-free shift}

\begin{proposition}[RH-free gauge fix, NEU-260a $\checkmark$]
\[
  \lambda_{\mathrm{can}}(a) := \lambda_a - 1, \qquad
  T_a^{\mathrm{can}} = A_a - (\lambda_a-1)I \ge I > 0.
\]
This is unconditional. The choice $\lambda=0$ requires RH
($A_a > 0$ holds under RH by Suzuki).
\end{proposition}

\section{Gauge freedom and intrinsic datum $U_a$}

\begin{proposition}[Gauge freedom, NEU-260b $\checkmark$]
Under $v_\pm \mapsto e^{i\alpha_\pm}v_\pm$, the parameter transforms
as $\theta\mapsto\theta+\beta-\alpha$. The intrinsic datum is the
unitary operator
\[
  U_a:\mathcal{N}_{+,a}\longrightarrow\mathcal{N}_{-,a},
\]
where $\mathcal{N}_{\pm,a} = \ker(\mathscr{D}_a^*\mp i)$.
The number $\theta$ is only a coordinate after a choice of bases.
\end{proposition}

\section{Arithmetic canonicity problem (open)}

\begin{conjecture}[Reflection symmetry, NEU-260b $?[O]$]
The Weil form satisfies
\[
  Q_W^a(f,f) = Q_W^a(\tilde f,\tilde f), \qquad
  \tilde f(x) := \overline{f(-x)}.
\]
If true, the reflection $x\mapsto -x$ provides the canonical
pairing $U_a^\sigma:\mathcal{N}_{+,a}\to\mathcal{N}_{-,a}$ with
$\theta_{\mathrm{can}}(a)=0$.
\end{conjecture}

\begin{conjecture}[Suzuki limit, Suzuki 2026]
\[
  e^{\phi(a,z)} W(a,\theta(a);z) \longrightarrow
  \frac{z^2 \xi(1/2-iz)}{\xi'(1/2-iz)},
  \quad a\to\infty,\text{ locally uniform.}
\]
Suzuki notes that $\phi\equiv 0$ may suffice.
\end{conjecture}

\section{Objekt-X hypothesis}

\begin{definition}[Objekt X, NEU-259/260 $?[O]$]
\[
  \text{Objekt X} := \left\{\mathcal{H}(T_{a,\lambda_{\mathrm{can}}(a)}),\;
  J_{a,b},\; \overline{\mathscr{D}}_{a,\theta(a)}\right\}_{0<a<b},
\]
with arithmetic-canonical $\theta(a)$, $\phi(a,z)$, $J_{a,b}$.
Under RH: $\mathcal{K}_X := \varinjlim_a\mathcal{H}(T_a)
\xrightarrow{\mathrm{RH}}\mathcal{H}_W\cong\ell^2(\Gamma,m_\gamma)$.
\end{definition}

\section{Open problems (P04)}

\begin{enumerate}
\item $Q_W^a$-reflection symmetry $\Rightarrow$ $\theta_{\mathrm{can}}(a)=0$? (NEU-260b)
\item Grenznormalisierung $\phi(a,z)$: arithmetic determination? (NEU-260c)
\item Canonical $J_{a,b}$ with intertwining
  $J_{a,b}\overline{\mathscr{D}}_{a,\theta(a)}\subset
  \overline{\mathscr{D}}_{b,\theta(b)}J_{a,b}$? (NEU-260d)
\end{enumerate}

\section{Source Nodes}

NEU-259 (final, Patch 2), NEU-260a (Patch), NEU-260b
(active/).

\end{document}
